The variance correction allowed fitted full-return generators to be
reused in a single-index model without adding market variance twice
(Fig.~\ref{fig:preservation}; Table~\ref{tab:aggregate}).
Its practical value was avoiding a separate generator fit for each
asset's regression residuals while achieving comparable distributional
fit. The corrected method passed $73.6\%$ of the full-length training
KS comparisons and retained heavy tails (Table~\ref{tab:aggregate}).
On the $2025$ holdout, it raised the mean cross-ticker KS pass rate
from $83.5\%$ under centered naive composition to $87.0\%$ and held
the median ratio of synthetic to observed variance at $0.97$
(Table~\ref{tab:oos_composition}). JumpHMM-on-residuals had a slightly
higher pass rate in training and a slightly lower rate in the holdout.
The comparison supported reuse as a practical option when full-return
generators are already available; fitting directly to residuals
remained a competitive alternative. Neither method was best for every
asset, and the average improvement concealed a subset of corrected
paths whose fit deteriorated relative to matched-length training
blocks.
The validation established closer marginal fit, but it did not establish
reliable tail coverage (Table~\ref{tab:oos_composition};
Table~\ref{tab:var}). The corrected method and JumpHMM-on-residuals
both had a one-day left-tail $99\%$ VaR exceedance rate of $1.67\%$,
above the nominal $1\%$ rate (Table~\ref{tab:var}). The $249$-day
holdout contained too few expected extreme events per asset for a
precise ranking. Distributional pass rates also require care because
the two-sample KS and AD tests assume independent observations,
whereas several evaluated generators produced serial dependence. We
therefore interpreted pass rates alongside distributional distances,
tail summaries, and autocorrelation errors. Matched-length training
blocks reduced the difference in test power between the training and
holdout comparisons, but the observed $2025$ history still represented
only one market period. The evidence supported improved variance and
distributional fit across the evaluated universe, with substantial
uncertainty about performance in other periods and for individual
assets.

The jump comparison showed that improving the single-asset generator
did not automatically improve the composed paths
(Table~\ref{tab:model_comparison}; Table~\ref{tab:jump_ablation}). On SPY, HMM-WJ
reduced absolute-return autocorrelation error relative to HMM-NJ while
retaining high distributional pass rates. Other generators achieved
lower autocorrelation error or higher pass rates, so the single-asset
result supported a tradeoff between temporal and marginal fit
(Table~\ref{tab:model_comparison}). With multi-asset composition,
market-only jumps improved training temporal fit with little change
in the KS pass rate, whereas adding asset-level episodes produced a
larger temporal improvement at a substantial marginal-fit cost.
Both settings increased autocorrelation error in $2025$, although
the variance correction remained effective
(Table~\ref{tab:jump_ablation}). These tests transferred one SPY
calibration from volume-weighted average prices to closing-price
models. Separate calibrations for NVDA, JNJ, and JPM had strong
in-sample fit, but three assets did not establish a general calibration
rule (Table~\ref{tab:cross_asset}). The asset episodes also came from
full-return models and were not identified separately as idiosyncratic
events. Per-asset tuning or distinct shared and asset-specific episode
models could change the tradeoff, but neither was tested here.

The correction's effect on the return distribution depended on the
asset's market exposure (Fig.~\ref{fig:preservation};
Table~\ref{tab:beta_bucket}). Corrected KS pass rates declined from
$87.6\%$ in the lowest $\beta$ quartile to $65.9\%$ in the highest
because the market path accounted for more of the composed variance
(Table~\ref{tab:beta_bucket}). Preserving total variance therefore
did not preserve the complete shape of the original return
distribution. The loading adjustment was unnecessary at observed
market variance but activated for most assets when market volatility
was doubled or tripled (Table~\ref{tab:stress}). Only two tracker
assets used the separate $R^2$-preserving branch
(Table~\ref{tab:branch}), so evidence for that branch relied mainly
on the synthetic-tracker checks. These properties also depended on
generating asset and market draws independently. If their covariance
were nonzero, the variance calculation would need an additional
cross-term. The correction thus addressed a specific composition
problem, with explicit assumptions about the supplied paths and
adjustments when the market term consumed too much of the target
variance.
Dependence among asset-specific residuals remained outside the
composition rule (Algorithm~\ref{alg:hybrid}; Online
Appendix~\ref{sec:supp-bias}). A shared market path linked same-day asset returns and
could produce simultaneous large moves, but independent residual
draws omitted the sector and style dependence present in real assets.
Omitting positive residual covariance can overstate the diversification
return from periodic rebalancing~\cite{boothFama1992,
boucheyNemtchinovWong2015, markowitz1976longRun}; Online
Appendix~\ref{sec:supp-bias} derives the contribution for this
construction. The limitation also applies to the residual-fit
baselines when their residuals are generated independently. Shifting
paths leaves cross-asset covariance unchanged, and rescaling
independent residuals cannot create the missing dependence. A
synchronized multivariate block bootstrap, residual factor model, or
copula could address that omission, while separate univariate models
would improve only within-asset behavior. The present evaluation
measured per-asset fit and VaR coverage, so it did not validate
portfolio-level risk or rebalancing performance. Those uses require a
joint residual model and direct evaluation of the resulting dependence.
